\section*{How this book works}

Every topic has four layers. Read them in this order; the first two take about two minutes.

\begin{enumerate}
\item \textsf{\color{sumcol}2-minute summary} -- the whole topic on half a page: the one idea, the
  one formula, the words the examiner wants to hear.
\item \textsf{\color{piccol}Picture it} -- an everyday image (an analogy). You remember a picture
  much longer than a sentence, and it gives you words when you are nervous.
\item \textsf{\color{brdcol}Board drill} -- the formula you must be able to write on the
  whiteboard. Write it by hand three times, without looking. Hands learn faster than eyes
  (tactile, or kinaesthetic, learning), and on the day you will stand at a board.
\item \textsf{\color{qcol}Questions} -- predicted exam questions. \textsf{Say} is what you say
  out loud, \textsf{Write} is what you put on the board, \textsf{Follow-up} is the next question
  that usually comes.
\end{enumerate}

\textsf{Who asks what} (from the list of past students):
\begin{itemize}
\item Prof.\ Klambauer: empirical risk, loss, gradient descent, backprop and the delta error,
  generalization. ``Write the simplest version first.''
\item Prof.\ Hochreiter: delta error, quick questions (regularization, optimizers, activations,
  the activation from Linz), REINFORCE, and everything LSTM.
\item Others (Rekabsaz, Lehner, Schl\"uter): Transformers and attention, RNNs, SVM, VC dimension,
  universal approximation.
\item Prof.\ Pirker (elective track): rCFD, your curves, CFD basics. Chapter 7.
\end{itemize}

\textsf{Study plan (suggested).} Day 1: Chapters 1--2 plus all board drills. Day 2: Chapters 3--4.
Day 3: Chapters 5--6. Day 4: Chapters 7--8 and the thesis. Then repeat only the summaries and
board drills every morning.
